\chapter{System Analysis and Requirements}

\section{Existing System}

In a typical small or medium-sized Kubernetes deployment, security is handled by a
combination of image scanning in the CI pipeline, Kubernetes RBAC and network
policies, and manual inspection of logs with \texttt{kubectl logs} or a log
aggregation stack. The first prototype of this project (``v0'') followed the same
spirit: a BCC probe captured only the file name of executed programs, an offline log
parser stored them in PostgreSQL and an Isolation Forest flagged unusual command names.

\subsection{Limitations of the Existing System}
\begin{itemize}
    \item \textbf{No runtime visibility:} scanners and policies do not see process, file
          and network behaviour after a pod starts; application logs rarely record an
          attacker's shell commands.
    \item \textbf{Narrow sensor:} the v0 probe captured only \texttt{execve} file names,
          not arguments, files, network destinations or privilege changes.
    \item \textbf{No Kubernetes context:} events could not be linked to a Pod, namespace
          or workload, so an operator could not tell which application was affected.
    \item \textbf{Weak detection:} the unsupervised Isolation Forest detected only 12\% of
          attacks on our test set (Chapter~8) and produced string-matching false positives,
          for example \texttt{runc} matched as netcat because it contains ``nc''.
    \item \textbf{No risk scoring or explanation:} there was no severity, no MITRE
          technique and no matched rule to explain an alert.
    \item \textbf{Train/serve mismatch:} the model was trained on full command lines while
          the probe delivered only binary paths.
    \item \textbf{No evaluation:} detection rate, false positives, latency and overhead
          were never measured.
\end{itemize}

\section{Proposed System}

The proposed system is a complete runtime IDS built on the pipeline
\emph{Observe $\rightarrow$ Enrich $\rightarrow$ Detect $\rightarrow$ Score $\rightarrow$
Alert $\rightarrow$ Evaluate}.

\subsection{Features of the Proposed System}
\begin{itemize}
    \item \textbf{Kernel-level sensing with eBPF:} a libbpf CO-RE program on 16
          system-call tracepoints captures process execution with arguments, sensitive
          file opens, network connections, \texttt{ptrace}, \texttt{setuid},
          \texttt{mount}, \texttt{unshare}, \texttt{setns}, \texttt{chmod}, kernel-module
          loading and \texttt{memfd\_create}.
    \item \textbf{Kubernetes enrichment:} every event is resolved from its cgroup to Pod,
          container, namespace, workload, image, service account and node.
    \item \textbf{Shared feature extraction:} one Python module computes 74 features for
          both training and live inference, removing train/serve skew.
    \item \textbf{Hybrid detection:} a 10-class Random Forest and 19 deterministic rules
          with exemptions for platform components.
    \item \textbf{Risk scoring and alerting:} 0--100 risk score, Low / Medium / High
          severity, MITRE ATT\&CK technique and de-duplication of repeated alerts.
    \item \textbf{Dashboard:} React web interface with overview, alerts, events and
          evaluation pages, protected by JWT login.
    \item \textbf{Built-in evaluation:} offline model metrics, live agent heartbeats
          (CPU, memory, latency, dropped events) and an automated attack simulation.
\end{itemize}

\subsection{Advantages of the Proposed System}
\begin{itemize}
    \item Observes what workloads actually do at runtime, without modifying them.
    \item Gives every alert an owner (Pod, namespace, workload) and an explanation (rules,
          ML probability, MITRE technique).
    \item Detects both known patterns (rules) and variants of them (ML).
    \item Reduces false positives on legitimate platform behaviour through exemptions and
          edge-case training data.
    \item Runs as a standard DaemonSet with low CPU and memory overhead.
    \item Reports its own accuracy and overhead, so its limits are visible.
\end{itemize}

\section{Feasibility Study}

\begin{itemize}
    \item \textbf{Technical feasibility:} eBPF CO-RE is supported by all mainstream Linux
          kernels since 5.8; libbpf, scikit-learn, FastAPI, MongoDB and React are mature
          open-source technologies.
    \item \textbf{Economic feasibility:} all components are free and open source and the
          whole system runs on a single laptop with minikube.
    \item \textbf{Operational feasibility:} the system is deployed with one script and used
          through a web browser; no change to monitored applications is required.
\end{itemize}

\section{Requirement Analysis}

\subsection{Functional Requirements}

\subsubsection*{1. Event Monitoring}
\begin{itemize}
    \item Capture process execution (path and arguments), sensitive file opens, network
          connections and security-sensitive system calls on every node.
    \item Filter uninteresting events inside the kernel to limit overhead.
    \item Report kernel-side drops (ring-buffer full) so data loss is visible.
\end{itemize}

\subsubsection*{2. Kubernetes Enrichment}
\begin{itemize}
    \item Map every event to Pod, container, namespace, workload, image, service account
          and node.
    \item Store events without Kubernetes context (host processes) as ``raw''.
\end{itemize}

\subsubsection*{3. Detection and Scoring}
\begin{itemize}
    \item Extract a fixed feature vector per event including a 60-second behaviour window.
    \item Classify each event with the ML model and evaluate the rule set.
    \item Produce a risk score (0--100), severity (Low / Medium / High), attack category
          and MITRE ATT\&CK technique.
    \item Merge identical alerts from the same container within ten minutes.
\end{itemize}

\subsubsection*{4. Storage and API}
\begin{itemize}
    \item Store events, alerts, agent metrics, evaluation results and users in MongoDB.
    \item Expire old events (24 h) and alerts (30 days) automatically.
    \item Provide REST endpoints for ingestion (agent token) and for the dashboard (JWT).
\end{itemize}

\subsubsection*{5. Dashboard}
\begin{itemize}
    \item Secure login for administrators.
    \item Overview of events, open alerts, severities, attack categories, riskiest pods,
          agent overhead and detection latency.
    \item Search and filter alerts and events; acknowledge or resolve alerts.
    \item Show offline and live evaluation results.
\end{itemize}

\subsubsection*{6. Evaluation}
\begin{itemize}
    \item Compute accuracy, precision, recall, F1-score, false-positive rate and ROC-AUC
          offline.
    \item Run harmless attack and benign scenarios in a test pod and measure detection
          rate, false positives, latency and execution overhead.
\end{itemize}

\subsection{Non-Functional Requirements}
\begin{itemize}
    \item \textbf{Performance:} detection latency below one second; agent CPU below 10\%
          of one core at normal event rates; zero dropped kernel events.
    \item \textbf{Security:} JWT authentication for users, a separate shared token for
          agents, bcrypt password hashing, random credentials in a Kubernetes Secret.
    \item \textbf{Reliability:} a failing rule or a malformed event must never stop the
          sensor.
    \item \textbf{Scalability:} one agent per node through a DaemonSet; stateless backend.
    \item \textbf{Usability:} explainable alerts and a simple web interface.
    \item \textbf{Maintainability:} modular code with unit tests; one feature module shared
          by training and inference.
    \item \textbf{Portability:} CO-RE eBPF binary runs on any kernel with BTF without
          recompilation.
\end{itemize}

\subsection{Hardware Requirements}

\begin{table}[H]
\centering
\caption{Hardware requirements}
\label{tab:hw}
\begin{tabular}{lll}
\toprule
\textbf{Component} & \textbf{Minimum} & \textbf{Recommended} \\
\midrule
Processor & 2 cores, x86\_64 & 4 cores \\
RAM & 6 GB & 8 GB or more \\
Free disk space & 15 GB & 20 GB \\
Network & Internet for first image build & -- \\
\bottomrule
\end{tabular}
\end{table}

\subsection{Software Requirements}

\begin{table}[H]
\centering
\caption{Software requirements}
\label{tab:sw}
\begin{tabular}{lll}
\toprule
\textbf{Software} & \textbf{Version (tested)} & \textbf{Purpose} \\
\midrule
Ubuntu Linux & 22.04 / 24.04, kernel 6.8 & Host OS, kernel $\geq$ 5.8 with BTF \\
Docker & 24+ & Image build, minikube driver \\
minikube & 1.32+ & Local Kubernetes cluster \\
kubectl & 1.28+ & Cluster management \\
Python & 3.12 & Agent, ML, backend, evaluation \\
Node.js + npm & 18+ & Dashboard development \\
clang / libbpf & (inside Docker build) & Compile eBPF program \\
Git / GitHub & any & Version control \\
\bottomrule
\end{tabular}
\end{table}
